-------------------------------------------------------------------------------\\
\textbf{Points to discuss:}
\begin{enumerate}
	\item{\textit{Argumentation Framework}}
	\item{\textit{Abstract argumentation semantics}}
	\item{\textit{Approaches to computing abstract argumentation semantics}}
	\item{\textit{Review of the solvers}}
\end{enumerate}
-------------------------------------------------------------------------------\\


Argumentation framework is central to the theory of abstract argumentation \cite{baroni2011introduction}. The argumentation framework is defined as a pair of a set of arguments, and a binary relation representing the attack relationship between arguments \cite{dung1995}. 

\theoremstyle{definition}
\begin{definition}{Argumentation Framework}
\label{AFdef}\\
An argumentation framework is a pair \textit{AF} = $<$\textit{AR, attacks}$>$ in which \textit{AR} is a set of finite arguments, and \textit{attacks} is a binary relation on \textit{AR}, hence \textit{attacks} $\subseteq$ \textit{AR} $\times$ \textit{AR}, where \textit{AR} $\times$ \textit{AR} = \{(\textit{a}, \textit{b}) $\vert$ \textit{a} $\in$ \textit{AR} and \textit{b} $\in$ \textit{AR}\}
\end{definition}

In definition \ref{AFdef}, \textit{AR} represents a set of arguments and \textit{attacks} represents set of pairs of arguments (\textit{a, b}), where (\textit{a, b}) $\in$ \textit{attacks}. Each pair of arguments in \textit{attacks} represents two arguments being in conflict. Hence, the arguments \textit{a} and \textit{b} are in conflict and the meaning of of \textit{attacks(a, b)} is that \textit{a} represents an attack against \textit{b} \cite{dung1995}. Hence, it can be said that the set of arguments \textit{AR} is conflict free if there are no arguments \textit{a} and \textit{b} in \textit{AR} such that \textit{A} attack \textit{B} \cite{dung1995}.
\newline
 
 
The argumentation framework can be represented as directed graph where the nodes represent abstract arguments and edges the attack relation. This can be seen in figure \ref{fig:argumentationFrameworkFigure}, where argument a attacks argument b, which in turn attacks argument c. C is also attacked by argument d.

\begin{figure}[h]
	\centering
	\includegraphics[width=\textwidth,keepaspectratio=true]{argumentationFramework.png}
	\caption{Representation fo argumentation framework, from \cite{argumentationFrameworkExample}}
	\label{fig:argumentationFrameworkFigure}
\end{figure}
 
Dung in his paper \cite{dung1995} has also defined notions of \textit{acceptable} and \textit{admissible} arguments, which are as follow:

\begin{definition}{Acceptable argument}
\label{AcceptableArgDef}\\
An argument \textit{a} $\in$ \textit{AR} is said to be \textit{acceptable} with respect to set \textit{S} $\subseteq$ \textit{AR} if and only if for each argument \textit{b} $\in$ \textit{AR} such that (\textit{b}, \textit{a}) $\in$ \textit{attacks}, there are some arguments \textit{c} $\in$ \textit{S} such that (\textit{c}, \textit{b}) $\in$ \textit{attacks}
\end{definition}

Based on the definition \ref{AcceptableArgDef}, definition of \textit{admissible} set can be concluded:

\begin{definition}{Admissible argument}
\label{AdmissibleArgDef}\\
A set of arguments \textit{S} is \textit{admissible} if and only if each argument in \textit{S} is acceptable with respect to \textit{S}.
\end{definition}

That means that the set of arguments can be described as admissible only if it is conflict-free. Which in turn means that no argument in S attack another argument in S.

Furthermore, Dung in his paper \textit{On the acceptability of arguments and its fundamental role in nonmonotonic reasoning, logic programming and n-person games} introduced the concept of semantics of abstract argumentation. The original concept included complete, stable, preferred and grounded semantics \cite{dung1995}, which was extended by ideal \cite{dung2007computing}, and semi-stable \cite{caminada2006semi} semantics.

%\paragraph{Preferred Semantics}\\

